\chapter{AI 编译器：从推理图到机器执行计划}

\section{推理引擎里为什么会出现编译器}

前一章把本地 7B 级推理引擎拆成模型加载、tokenizer、decoder graph、KV cache、CPU/GPU 后端和证据报告。这里要再向下压一层：真实引擎不能每次遇到一个算子就临时解释执行。它必须把模型结构、张量形状、量化格式、设备能力和运行时状态编成一个可执行计划。这个过程就是 AI 编译器在推理引擎中的位置。

不要把 AI 编译器理解成“把 Python 变成机器码”的窄概念。推理场景里，输入通常已经不是普通源代码，而是模型文件、配置、tokenizer 规则、权重量化元数据和一张 decoder-only 计算图。编译器要做的事情是把这张图降到后端能高效执行的 kernel 序列，同时保持 logits、KV cache 和采样状态的语义不变。

\begin{lstlisting}[numbers=none,caption={推理引擎中的 AI 编译流水线}]
model/config/tokenizer
  -> graph import and shape checking
  -> typed tensor IR
  -> graph rewrite and operator fusion
  -> memory planning and lifetime analysis
  -> backend lowering
  -> kernel selection and packed layout generation
  -> executable plan
  -> runtime dispatch for prefill/decode
\end{lstlisting}

这条流水线和传统编译器很像。模型文件相当于源程序，图导入相当于 parsing，shape/dtype/quantization 检查相当于 semantic analysis，typed tensor IR 相当于中间表示，fusion 和 layout 选择相当于优化，CPU/GPU lowering 相当于后端代码生成，executable plan 则相当于目标程序。区别在于它编译的对象不是标量控制流为主的 C++ 函数，而是张量、状态、量化参数和设备 kernel 组成的推理程序。

\begin{keyidea}
AI 编译器在推理引擎中的核心职责，是把稳定的模型语义编成稳定的机器执行计划。它既不是单个 kernel，也不是普通 runtime 调度器，而是连接图语义、张量布局、内存生命周期、后端能力和正确性证据的编译流水线。
\end{keyidea}

\section{前端：导入图时先建立语义合同}

普通编译器前端会把源代码变成 token、AST 和符号表。AI 编译器的前端面对的是模型结构。以 decoder-only Transformer 为例，前端至少要确认这些事实：

\begin{itemize}
  \item 每层有哪些算子：RMSNorm、Q/K/V projection、RoPE、attention、Wo projection、MLP gate/up/down、残差。
  \item 每个张量的 dtype、shape、stride、量化参数和驻留位置。
  \item 每个权重是否已经按某种 group size、scale 类型和物理布局保存。
  \item prefill 和 decode 是否共享同一张图，还是需要两套 specialization。
  \item KV cache 的逻辑 shape、最大上下文、当前位置和写入规则。
\end{itemize}

这些检查不是文档整理，而是编译器不变量。若 \code{Wq} 的输入维度和 hidden size 对不上，后端 kernel 不应该继续猜；若某个 int4 权重缺少 group scale，量化 lowering 不能把它当 int8；若 KV cache 的 head 数和当前 layer 的 attention head 数不一致，decode 会在很晚的位置产生错误 logits。前端要么拒绝模型，要么产出满足合同的 typed tensor IR。

\begin{lstlisting}[numbers=none,caption={typed tensor IR 的最小信息}]
%0 = input_tokens       : tensor<[T], i32>
%1 = embedding(%0)      : tensor<[T,H], f16>
%2 = rmsnorm(%1, gamma) : tensor<[T,H], f16>
%3 = linear(%2, Wq_i4)  : tensor<[T,Q], f16>
%4 = linear(%2, Wk_i4)  : tensor<[T,K], f16>
%5 = linear(%2, Wv_i4)  : tensor<[T,V], f16>
%6q, %6k = rope(%3, %4, pos) : tensor<[T,Q], f16>, tensor<[T,K], f16>
%7 = kv_append(%6k, %5)      : state<KVCache>
%8 = attention(%6q, %7)      : tensor<[T,H], f16>
\end{lstlisting}

这段 IR 不是为了好看。它让每条边都带上类型、形状和状态含义。没有这些信息，后面的融合、内存复用、layout 转换和后端选择都只能靠脆弱约定。

\section{IR：图 IR、张量 IR 和循环 IR 要分层}

AI 编译器最容易犯的错误，是用一层表示承载所有事情。图 IR 擅长表达算子依赖：RMSNorm 后接 Q/K/V projection，RoPE 后写 KV cache，attention 后接 Wo projection。张量 IR 擅长表达 shape、stride、dtype、broadcast、reduction 和 layout。循环 IR 或 kernel IR 才适合表达 tile、向量宽度、展开、线程划分和寄存器累加器。

\begin{systemdiagram}{AI 编译器的三层 IR}
\begin{pdfdiagram}[x=1cm,y=1cm]
  \node[book accent node,text width=2.35cm] (graph) at (0,1.35) {Graph IR\\算子依赖};
  \node[book teal node,text width=2.35cm] (tensor) at (3.15,1.35) {Tensor IR\\shape/layout};
  \node[book purple node,text width=2.35cm] (loop) at (6.3,1.35) {Loop IR\\tile/vector};
  \node[book amber node,text width=2.35cm] (target) at (9.45,1.35) {Target plan\\CPU/GPU};

  \node[book label,text width=2.35cm] at (0,-0.25) {fusion\\state edges\\dependency};
  \node[book label,text width=2.35cm] at (3.15,-0.25) {dtype\\stride\\quant};
  \node[book label,text width=2.35cm] at (6.3,-0.25) {unroll\\SIMD\\threads};
  \node[book label,text width=2.35cm] at (9.45,-0.25) {kernel\\layout\\sync};

  \draw[book strong arrow] (graph) -- (tensor);
  \draw[book strong arrow] (tensor) -- (loop);
  \draw[book strong arrow] (loop) -- (target);
\end{pdfdiagram}
\diagramnote{本图要证明的命题是：AI 编译器不能把算子图、张量布局和机器循环混成一层，否则优化会互相污染。}
\end{systemdiagram}

分层的好处是边界清楚。Graph IR 可以做无关后端的代数改写，例如把连续的 reshape 消掉，或把 RMSNorm 与后续线性层之间的临时 buffer 标成可融合候选。Tensor IR 可以决定某个 \code{[tokens, hidden]} 张量是否连续、某个 int4 权重是否要按 group 解码、某个 broadcast 是否会制造隐式 stride 0。Loop IR 再决定 CPU 上一个 tile 处理多少列、用几个累加器、尾部怎样 mask，或者 GPU 上一个 block 负责多少 token 和 head。

传统编译器也有类似分层。AST 不应该直接塞寄存器分配细节；SSA IR 不应该直接保存源代码排版；机器 IR 才关心调用约定和寄存器。AI 编译器同理：高层 graph 不应该知道 AVX2 的寄存器数，低层 kernel 也不应该重新推断模型层数和 KV cache 语义。

\section{优化：融合不是越多越好}

算子融合是 AI 编译器最常见的优化，但它经常被误讲成“把算子合在一起就快”。实际判断必须回到内存流和生命周期。若两个算子之间的中间张量很大、只被消费一次、且融合后不会增加过多寄存器压力，那么融合可以减少一次写回和一次读取。若中间结果还被 residual、debug oracle、另一个分支或 KV cache 使用，盲目融合会破坏语义或让 buffer 生命周期变复杂。

\begin{lstlisting}[numbers=none,caption={融合决策要检查的事实}]
candidate: RMSNorm -> Linear

check:
  output of RMSNorm has exactly one consumer
  no oracle boundary requires materialized RMSNorm output
  fused kernel can stream input and gamma without spilling too much
  quantized Linear still has access to scale/group metadata
  backend supports the fused dtype/layout combination

if all true:
  lower to fused_rmsnorm_linear
else:
  keep two operators and expose the materialized buffer
\end{lstlisting}

attention 更能说明问题。prefill attention 可能是大块矩阵计算，decode attention 是一个 query 读取越来越长的 K/V cache。它们的形状不同，适合的 kernel 不同，融合边界也不同。prefill 可以倾向于高吞吐 tile；decode 更在意小 batch latency、KV cache 读流和 softmax 的数值稳定。一个“通用 attention 融合”若不区分这两种模式，很容易在其中一种场景下退化。

\begin{warningbox}
融合优化必须由数据依赖、生命周期、数值误差和后端能力共同决定。只看算子名字相邻，就把它们写进一个大 kernel，是把编译器优化退化成手工拼接。
\end{warningbox}

\section{内存规划：生命周期分析是推理编译器的核心}

本地推理最贵的不只是算术，还有内存容量和带宽。权重、packed weight、activation、KV cache、临时 workspace 和输出 logits 的生命周期完全不同。AI 编译器要像传统编译器做寄存器分配一样，为张量做内存分配。

\begin{lstlisting}[numbers=none,caption={张量生命周期的简化账本}]
weights:
  live from model load to session destroy
  read-only, often mmap or packed once

packed weights:
  live after backend preparation
  backend-specific physical layout

activation:
  live inside one layer or a small operator group
  can often reuse arena slots

KV cache:
  live across all decode steps
  grows by position and cannot be reused as scratch

workspace:
  live during one kernel or one scheduled segment
  size depends on backend algorithm
\end{lstlisting}

这就是为什么内存规划不能靠 \code{new}/\code{delete} 随手分配。activation 可以进入 arena 并按生命周期复用；KV cache 必须按最大上下文或可增长策略预留；packed weight 可能比原始权重多占一份内存，但换来热路径布局稳定；GPU workspace 要考虑 stream 同步和设备内存峰值。若编译器没有显式生命周期，runtime 只能在执行时保守分配，结果是峰值内存上升、cache locality 变差、调试也更困难。

这一步和传统编译器的 liveness analysis 很接近。变量最后一次使用后可以释放寄存器；张量最后一次消费后可以释放 arena slot。区别在于张量尺寸大得多，KV cache 还跨 token 存活，后端 layout 可能要求额外对齐和 padding。

\section{后端 lowering：CPU 和 GPU 的目标不同}

同一个 Tensor IR 降到 CPU 和 GPU 时，目标函数不同。CPU 后端关心 cache line、TLB、SIMD 宽度、指令吞吐、分支、线程划分、NUMA 和 false sharing。GPU 后端关心 device memory、coalesced load、shared memory、warp/block 划分、kernel launch、stream 同步和 tensor core 格式。AI 编译器的后端 lowering 不能假装它们只是两个函数指针。

\begin{lstlisting}[numbers=none,caption={同一 linear 算子的两个 lowering 方向}]
linear(input[T,H], weight[O,H]):

CPU decode lowering:
  choose GEMV or small-batch GEMM
  use packed weight panels
  select SIMD width and tail handling
  partition output rows across worker threads

GPU prefill lowering:
  choose batched GEMM kernel
  ensure device-resident input and weight
  select vendor layout or internal tile layout
  schedule stream dependencies
\end{lstlisting}

这里的关键是 specialization。prefill 的 \code{T} 可能较大，decode 的 \code{T} 通常是 1；长上下文 attention 和短上下文 attention 也不该强行共享同一个计划。编译器可以为不同 shape、batch、上下文上限和后端生成多个 executable plan，再让 runtime 在请求进入时选择。这样做比每次动态判断更复杂，但能把热路径从“解释图”变成“执行计划”。

\section{运行时：编译计划仍然需要调度}

编译器生成计划后，runtime 仍然有工作。它要管理 session，接收 prompt，维护当前 position，选择 prefill 或 decode 计划，推进 KV cache，处理 sampler，记录 profiler 事件，并在必要时切换后端或回退到 reference。编译器负责把能提前决定的事情提前决定，runtime 负责处理请求期才知道的状态。

\begin{lstlisting}[numbers=none,caption={编译期和运行期的责任划分}]
compile time:
  validate model
  build typed graph
  choose layouts and fusion
  allocate static arenas
  prepare packed weights
  build backend executable plans

run time:
  tokenize prompt
  choose prefill/decode plan by shape and backend
  bind session buffers and KV position
  dispatch kernels
  sample next token
  collect correctness and performance evidence
\end{lstlisting}

这个划分也解释了为什么“推理引擎”和“AI 编译器”不是二选一。没有编译器，引擎会在每次请求中重复解释图、判断 layout、临时选择 kernel；没有 runtime，编译器生成的计划没有 session、KV cache 和 sampler 可以驱动。成熟系统通常是二者配合：前者把模型变成计划，后者把计划用在一个个真实请求上。

\section{runtime 调度器：执行计划怎样变成 token 流}

编译器产出 executable plan 后，runtime 调度器要把它用在真实请求上。调度器至少处理五件事：请求进入、prefill 排队、decode 循环、取消/超时、资源释放。它不是简单循环调用 \code{plan.run()}。

\begin{lstlisting}[numbers=none,caption={runtime 调度器的最小状态}]
Scheduler:
  loaded_models
  plan_cache
  ready_prefill_queue
  active_decode_sessions
  memory_reservations
  backend_contexts
  trace_sink
\end{lstlisting}

prefill 和 decode 队列要分开。prefill 可能是长 prompt 的大块计算，decode 是很多 session 的单 token 步进。若长 prefill 一直占住所有工作线程或 GPU stream，已经在交互中的请求会断流；若 decode 过度优先，长 prompt 首 token 会等太久。调度策略要根据业务目标取舍。

\begin{center}
\begin{tabular}{p{0.20\linewidth}p{0.34\linewidth}p{0.30\linewidth}}
\toprule
策略 & 优点 & 风险 \\
\midrule
prefill 优先 & 首 token 更快进入模型 & 活跃会话 decode 抖动 \\
decode 优先 & 流式输出更平滑 & 新请求 TTFT 上升 \\
分配时间片 & 两者较平衡 & 实现复杂，profile 更难读 \\
按 deadline & 符合 SLO & 需要更准的耗时估计 \\
\bottomrule
\end{tabular}
\end{center}

\section{取消、超时和资源释放}

业务服务必须支持取消。用户关闭连接、上游超时、内容策略中止、队列过载，都可能要求停止某个 session。取消不是只设置一个 bool。runtime 要保证：

\begin{itemize}
  \item 不再为该 session 提交新的 decode step。
  \item 已提交到 GPU stream 的工作要么等待完成再释放 buffer，要么走后端支持的安全取消路径。
  \item KV cache page、activation arena、workspace reservation 和 trace buffer 被释放。
  \item sampler state 和 token buffer 不再被其他线程访问。
  \item FinalResponse 标出 \code{cancelled} 或 \code{timeout}，方便业务层区分。
\end{itemize}

\begin{lstlisting}[numbers=none,caption={取消路径的资源顺序}]
cancel(session):
  mark session as cancelling
  stop scheduling new decode steps
  wait or fence backend work using this session buffers
  release KV pages and arena reservations
  close stream with final cancelled event
  write trace summary
\end{lstlisting}

这个顺序和第二册讲过的异步生命周期一致：释放资源之前，必须证明没有异步执行还会读写它。CPU 后端也有类似问题，只是表现为线程池任务；GPU 后端表现为 stream/event。

\section{plan cache 和 shape bucket}

如果每个请求都重新编译 plan，热路径会很慢。runtime 通常需要 plan cache。cache key 至少包含模型、backend、量化 profile、stage、shape bucket 和 debug/oracle 配置。

\begin{lstlisting}[numbers=none,caption={PlanCacheKey 字段}]
PlanCacheKey:
  model_id
  manifest_hash
  backend
  quant_profile_id
  stage: prefill | decode
  prompt_length_bucket
  context_length_bucket
  oracle_policy
  trace_level
\end{lstlisting}

shape bucket 是折中。完全为每个 prompt length 编译一个计划，cache 会爆；所有长度共用一个计划，kernel 可能低效。可以把 prompt length 分成 \code{1-32}、\code{33-128}、\code{129-512}、\code{513+} 等桶，把 decode context 也按 \code{0-512}、\code{512-2048}、\code{2048+} 分桶。每个桶对应不同 attention plan、workspace 和 profiler baseline。

debug/oracle 配置也要进入 key。因为打开 layer oracle 可能强制 materialize 某些中间张量，导致 fusion 和 memory planning 不同。不能把 debug plan 和 production plan 混用。

\section{正确性：优化后的计划必须能被验证}

AI 编译器的优化会改变执行顺序、buffer 物化点、累加精度、量化缩放和后端 kernel。它不一定能做到 bitwise identical，但必须有可解释的误差合同。对一个本地推理引擎来说，至少要比较这些层次：

\begin{itemize}
  \item 单算子输出：RMSNorm、linear、RoPE、attention、MLP 与 reference 的误差。
  \item 层输出：一个 decoder block 前后 hidden state 的误差。
  \item logits：最大绝对误差、top-k 排名变化和采样前分布变化。
  \item 生成序列：固定 seed 和固定 sampler 下 token 序列是否在合同内。
  \item 性能证据：每个优化计划的 prefill/decode 分项耗时和峰值内存。
\end{itemize}

传统编译器有 verifier、IR dump、debug info 和测试套件；AI 编译器也需要对应工具。Graph IR dump 用来查融合是否正确，memory plan dump 用来查 buffer 生命周期，backend plan dump 用来查 kernel 选择，oracle report 用来查误差，profiler report 用来查优化是否真的改善瓶颈。

\section{本章验收线}

读完本章，应该能把“AI 编译器”和“推理引擎”的关系讲清楚：

\begin{itemize}
  \item AI 编译器把模型、图、张量合同和后端能力编成 executable plan。
  \item 推理引擎在运行期维护 prompt、session、KV cache、sampler 和 profiler。
  \item Graph IR、Tensor IR 和 Loop IR 分别负责算子依赖、张量布局和机器循环。
  \item 融合优化必须检查消费者、生命周期、误差边界和后端支持。
  \item 内存规划本质上是张量级生命周期分析，KV cache 不能当普通 scratch buffer。
  \item CPU/GPU lowering 共享语义合同，但按不同硬件目标选择不同物理布局和 kernel。
\end{itemize}

后续进入真实模型加载、tokenizer 和具体算子时，这套编译视角会反复出现：每个优化都要说明它在哪一层 IR 生效，改变了哪条地址流或生命周期，后端计划如何验证，以及 runtime 怎样在 prefill 和 decode 中使用它。
